\begin{center}
{\Huge \scshape Qiong Liu, Ph.D.} \\ \vspace{4pt}
\small
\faPhone\ 475-280-1364 \quad
\href{mailto:qiongliu.work@gmail.com}{qiongliu.work@gmail.com} \quad
\href{https://www.linkedin.com/in/qiong-l-529b23196}{LinkedIn} \quad
\href{https://scholar.google.com/citations?user=SvnpHCkAAAAJ&hl=en}{Google Scholar}

Authorized to work in the U.S. without sponsorship
\end{center}

\section{Professional Summary}
\textbf{Healthcare AI / Medical Imaging Scientist} with expertise in machine learning for image reconstruction, denoising, motion correction, and quantitative analysis in PET imaging. Strong track record of developing clinically meaningful AI methods for low-count, dynamic, and motion-corrupted imaging data. Experienced across research and product-oriented environments, with first-author publications, patent-pending work, and end-to-end algorithm development experience.

\section{Core Expertise}
\textbf{Medical Imaging AI:} PET reconstruction, denoising, motion correction, quantitative imaging, parametric imaging, dynamic imaging \\
\textbf{Machine Learning:} Deep learning, self-supervised learning, diffusion models, transformers, multimodal learning, physics-informed AI \\
\textbf{Programming:} Python, PyTorch, NumPy, SciPy, MATLAB \\
\textbf{Strengths:} Algorithm development, quantitative evaluation, clinically informed modeling, cross-functional collaboration

\section{Experience}
\textbf{Canon Medical Research USA} \hfill Apr 2025 -- Present \\
\textit{AI Scientist (Reconstruction), Vernon Hills, IL}
\begin{itemize}[leftmargin=*]
\item Developing AI-driven motion correction methods for PET reconstruction to improve image quality and quantitative accuracy in dynamic and low-count imaging
\item Designing next-generation gating algorithms for respiratory and cardiac motion compensation in dynamic PET
\item Building reconstruction pipelines that integrate deep learning and physics-based modeling for clinically robust imaging
\item Collaborating across research and product teams to evaluate algorithm performance and translational potential
\end{itemize}

\textbf{Canon Medical Research USA} \hfill Jun 2023 -- May 2024 \\
\textit{Research Scientist Intern, Vernon Hills, IL}
\begin{itemize}[leftmargin=*]
\item Developed AI-based PET reconstruction and denoising methods for low-count imaging scenarios
\item Designed evaluation frameworks to assess image quality and quantitative consistency across datasets
\item Presented technical findings to guide product and algorithm development
\end{itemize}

\textbf{United Imaging Healthcare} \hfill Jul 2019 -- Aug 2019 \\
\textit{Student Intern, Wuhan, China}
\begin{itemize}[leftmargin=*]
\item Supported development and validation of engineering systems for clinical applications
\end{itemize}

\section{Patents}
\textbf{Edge-Guided Deformation Field Fusion for PET Motion Correction} \hfill 2026 \\
Patent Pending, \textit{Lead Inventor}
\begin{itemize}[leftmargin=*]
\item Proposed a spatially adaptive AI framework to improve lesion localization and quantitative accuracy in PET motion correction
\end{itemize}

\textbf{Automatic Data-Driven Cardiac-Respiratory Gating} \hfill 2025 \\
Patent Pending, \textit{Lead Inventor}
\begin{itemize}[leftmargin=*]
\item Developed an AI-based method for separating respiratory and cardiac motion signals from dynamic PET data
\end{itemize}

\section{Selected Projects}

\textbf{Physiology-Guided Motion Correction for Dynamic PET}
\begin{itemize}[leftmargin=*]
\item Developed AI-based motion correction methods for dynamic PET, addressing real-world challenges from heterogeneous tracer uptake, low-count noise, and anatomically ambiguous regions
\item Proposed a \textbf{physiology-guided spatial mask} to provide clinically meaningful guidance during network training, improving robustness in regions where conventional learning-based constraints were insufficient
\item Introduced \textbf{spatially adaptive regularization} in place of global smoothness constraints, enabling more reliable deformation estimation across heterogeneous image regions
\item Improved motion correction robustness and temporal consistency for downstream quantitative analysis in dynamic PET imaging
\end{itemize}

\textbf{Kinetics-Guided Dynamic PET Denoising}
\begin{itemize}[leftmargin=*]
\item Developed \textbf{Patlak-guided self-supervised learning} methods for dynamic PET denoising to improve the quality and quantitative reliability of parametric imaging
\item Incorporated tracer kinetic constraints directly into AI model training, preserving physiologically meaningful temporal behavior across dynamic frames
\item Demonstrated that integrating \textbf{physics- and kinetics-informed priors} into deep learning improves performance in noisy and limited-data clinical imaging settings
\item Improved downstream Patlak-based quantification and parametric consistency in whole-body dynamic \textsuperscript{18}F-FDG PET studies
\end{itemize}

\textbf{Interpretable Deep Learning for Low-Count and Dynamic PET}
\begin{itemize}[leftmargin=*]
\item Developed personalized, population-based, and deep image prior methods for PET denoising with emphasis on preserving clinically relevant quantitative information
\item Studied how training data composition and population priors affect denoising outcome and downstream quantification performance across heterogeneous imaging datasets
\item Analyzed CNN training behavior and evolving noise patterns to better understand model behavior in low-count and dynamic PET reconstruction tasks
\item Showed that deep learning for healthcare imaging, while often treated as a black box, can be systematically interpreted and optimized for more reliable clinical use
\end{itemize}

\section{Education}
\textbf{Yale University} \hfill Aug 2020 -- May 2025 \\
Ph.D. in Biomedical Engineering \\
Thesis: Improving Image Quality and Quantification Accuracy for Static and Dynamic PET

\textbf{Huazhong University of Science and Technology} \hfill Sep 2016 -- Jun 2020 \\
B.S. in Biomedical Engineering, Minor in Computer Science \\
GPA: 3.94/4.00

\section{Selected Publications}
\textbf{First Author}
\begin{itemize}[leftmargin=*]
\item Liu Q., et al., \textit{Parametric cardiac imaging with $^{18}$F-flutemetamol PET to evaluate the impact of tafamidis in patients with transthyretin cardiac amyloidosis}, \textbf{Journal of Nuclear Medicine}, 2026
\item Liu Q., et al., \textit{Patlak-guided self-supervised learning for dynamic PET denoising}, \textbf{IEEE Transactions on Radiation and Plasma Medical Sciences}, 2026
\item Liu Q., et al., \textit{Population-based deep image prior for dynamic PET denoising: a data-driven approach to improve parametric quantification}, \textbf{Medical Image Analysis}, 2024
\item Liu Q., et al., \textit{A personalized deep learning denoising strategy for low-count PET images}, \textbf{Physics in Medicine \& Biology}, 2022
\end{itemize}

\section{Awards}
\textbf{Society of Nuclear Medicine and Molecular Imaging (SNMMI)}
\begin{itemize}[leftmargin=*]
\item Young Investigator Award (1st Place), 2025
\item Poster Award (2nd Place), 2024
\item Young Investigator Award (2nd Place), 2023
\end{itemize}

\textbf{Mitacs}
\begin{itemize}[leftmargin=*]
\item Globalink Research Award, 2019
\end{itemize}